% ch0-preliminaries.tex — Chapter 0: Preliminaries
% Review of sets, mappings, relations, algebraic structures, and proof techniques.

\chapter{Preliminaries}\label{ch:preliminaries}

Linear algebra does not exist in a vacuum. It rests upon a small but
essential toolkit from set theory, logic, and abstract algebra. The purpose
of this chapter is to establish that toolkit clearly and concisely so that
the rest of the course may proceed without interruption.

Most of the material here should be familiar from earlier courses. We
therefore adopt a brisk pace: definitions are stated precisely, key
examples are worked through, and proofs are included only when they
illuminate a technique that will recur later. The reader who is
comfortable with all five sections may safely skip ahead to

ef{ch:vector-spaces}; the reader who is not should work through the
exercises at the end before continuing.


% ============================================================
\section{Sets}\label{sec:sets}
% ============================================================

\begin{definition}{Set}{def:set}
A \emph{set}\index{set} is a collection of distinct objects, called its
\emph{elements}. We write $x\in A$ to mean that $x$ is an element of~$A$,
and $x\notin A$ otherwise.
\end{definition}

We recall the standard number sets that will pervade every chapter:
\[
  \N \;\subset\; \Z \;\subset\; \Q \;\subset\; \R \;\subset\; \C.
\]
Here $\N = \set{0,1,2,\ldots}$ includes zero by convention.

\begin{definition}{Subset and set equality}{def:subset}
Let $A$ and $B$ be sets.
\begin{enumerate}[label=(\roman*)]
  \item $A$ is a \emph{subset}\index{subset} of $B$, written $A\subseteq B$,
    if every element of $A$ belongs to $B$.
  \item $A = B$ if and only if $A\subseteq B$ and $B\subseteq A$.
  \item $A$ is a \emph{proper subset}\index{proper subset} of $B$, written
    $A\subsetneq B$, if $A\subseteq B$ and $A\neq B$.
\end{enumerate}
\end{definition}

\begin{definition}{Set operations}{def:set-ops}
Let $A$ and $B$ be subsets of a universal set~$U$.
\begin{enumerate}[label=(\roman*)]
  \item \textbf{Union:} $A\cup B \deff \setcond{x\in U}{x\in A \text{ or } x\in B}$.
  \item \textbf{Intersection:} $A\cap B \deff \setcond{x\in U}{x\in A \text{ and } x\in B}$.
  \item \textbf{Difference:} $A\setminus B \deff \setcond{x\in U}{x\in A \text{ and } x\notin B}$.
  \item \textbf{Complement:} $A^c \deff U\setminus A$.
  \item \textbf{Cartesian product:} $A\times B \deff \setcond{(a,b)}{a\in A,\; b\in B}$.
  \item \textbf{Power set:} $\mathcal{P}(A) \deff \setcond{S}{S\subseteq A}$.
\end{enumerate}
\end{definition}

\begin{example}{Cartesian products}{ex:cartesian}
$\R^2 = \R\times\R$ is the Euclidean plane. More generally,
$\R^n = \underbrace{\R\times\cdots\times\R}_{n}$ is the set of all
$n$-tuples of real numbers, the primary setting for Chapters~1--3.
\end{example}


% ============================================================
\section{Mappings (functions)}\label{sec:mappings}
% ============================================================

\begin{definition}{Mapping}{def:mapping}
A \emph{mapping}\index{mapping} (or \emph{function}) from a set $A$ to a
set $B$ is a rule $f\colon A\to B$ that assigns to each element $a\in A$
exactly one element $f(a)\in B$. The set $A$ is the \emph{domain} and $B$
is the \emph{codomain}.
\end{definition}

\begin{definition}{Image and preimage}{def:image-preimage}
Let $f\colon A\to B$ be a mapping.
\begin{enumerate}[label=(\roman*)]
  \item The \emph{image}\index{image} of a subset $S\subseteq A$ is
    $f(S) \deff \setcond{f(a)}{a\in S}$.
  \item The \emph{image} (or \emph{range}) of $f$ is $\im f \deff f(A)$.
  \item The \emph{preimage}\index{preimage} of a subset $T\subseteq B$ is
    $f^{-1}(T) \deff \setcond{a\in A}{f(a)\in T}$.
\end{enumerate}
\end{definition}

\begin{definition}{Injectivity, surjectivity, bijectivity}{def:inj-surj-bij}
Let $f\colon A\to B$.
\begin{enumerate}[label=(\roman*)]
  \item $f$ is \emph{injective}\index{injective} (one-to-one) if
    $f(a_1)=f(a_2)$ implies $a_1=a_2$.
  \item $f$ is \emph{surjective}\index{surjective} (onto) if
    $\im f = B$, i.e.\ for every $b\in B$ there exists $a\in A$ with $f(a)=b$.
  \item $f$ is \emph{bijective}\index{bijective} if it is both injective
    and surjective. In this case $f$ has an \emph{inverse}
    $\inv{f}\colon B\to A$ satisfying $\inv{f}\circ f = \Id_A$ and
    $f\circ\inv{f} = \Id_B$.
\end{enumerate}
\end{definition}

\begin{example}{A classical injection and surjection}{ex:inj-surj}
\begin{enumerate}[label=(\alph*)]
  \item The map $f\colon\R\to\R$, $f(x)=x^2$, is neither injective
    (since $f(-1)=f(1)=1$) nor surjective (since $-1\notin\im f$).
  \item Restricting the domain and codomain: the map
    $g\colon [0,\infty)\to [0,\infty)$, $g(x)=x^2$, is bijective, with
    inverse $\inv{g}(y)=\sqrt{y}$.
\end{enumerate}
\end{example}

\begin{definition}{Composition}{def:composition}
If $f\colon A\to B$ and $g\colon B\to C$, their \emph{composition}\index{composition}
is the mapping $g\circ f\colon A\to C$ defined by $(g\circ f)(a) \deff g(f(a))$.
Composition is associative: $h\circ(g\circ f)=(h\circ g)\circ f$.
\end{definition}

\begin{proposition}{Composition preserves injectivity and surjectivity}{prop:comp-inj-surj}
Let $f\colon A\to B$ and $g\colon B\to C$.
\begin{enumerate}[label=(\roman*)]
  \item If $f$ and $g$ are injective, then $g\circ f$ is injective.
  \item If $f$ and $g$ are surjective, then $g\circ f$ is surjective.
  \item If $g\circ f$ is injective, then $f$ is injective.
  \item If $g\circ f$ is surjective, then $g$ is surjective.
\end{enumerate}
\end{proposition}

\begin{proof}
We prove (i); the others are similar exercises (see

ef{exo:exo:comp-inj-surj}).

Suppose $f$ and $g$ are injective and $(g\circ f)(a_1) = (g\circ f)(a_2)$.
Then $g(f(a_1)) = g(f(a_2))$. Since $g$ is injective,
$f(a_1) = f(a_2)$. Since $f$ is injective, $a_1 = a_2$.
\end{proof}


% ============================================================
\section{Relations}\label{sec:relations}
% ============================================================

\begin{definition}{Binary relation}{def:relation}
A \emph{binary relation}\index{relation!binary} on a set $A$ is a subset
$\mathcal{R}\subseteq A\times A$. We write $a\;\mathcal{R}\;b$ to mean
$(a,b)\in\mathcal{R}$.
\end{definition}

\subsection{Equivalence relations}\label{subsec:equiv}

\begin{definition}{Equivalence relation}{def:equiv-rel}
A binary relation $\sim$ on $A$ is an \emph{equivalence
relation}\index{relation!equivalence} if it is:
\begin{enumerate}[label=(\roman*)]
  \item \textbf{Reflexive:} $a\sim a$ for all $a\in A$.
  \item \textbf{Symmetric:} $a\sim b$ implies $b\sim a$.
  \item \textbf{Transitive:} $a\sim b$ and $b\sim c$ imply $a\sim c$.
\end{enumerate}
\end{definition}

\begin{definition}{Equivalence class and quotient set}{def:equiv-class}
Let $\sim$ be an equivalence relation on~$A$. The \emph{equivalence
class}\index{equivalence class} of $a\in A$ is
\[
  [a] \deff \setcond{b\in A}{b\sim a}.
\]
The \emph{quotient set}\index{quotient set} is
$A/{\sim} \;\deff\; \setcond{[a]}{a\in A}$.
\end{definition}

\begin{example}{Congruence modulo $n$}{ex:congruence}
Fix an integer $n\geq 1$. On $\Z$, define $a\sim b$ if and only if
$n\mid (a-b)$ (i.e.\ $a\equiv b\pmod{n}$). This is an equivalence
relation with $n$ equivalence classes: $[0],[1],\ldots,[n-1]$. The
quotient set is denoted $\Z/n\Z$ (or simply $\Z_n$).

For $n=3$: $[0]=\set{\ldots,-6,-3,0,3,6,\ldots}$,
$[1]=\set{\ldots,-5,-2,1,4,7,\ldots}$,
$[2]=\set{\ldots,-4,-1,2,5,8,\ldots}$.
\end{example}

\begin{remark}{Equivalence classes partition the set}{rem:partition}
The equivalence classes of any equivalence relation on $A$ form a
\emph{partition} of $A$: they are nonempty, pairwise disjoint, and their
union is~$A$. Conversely, every partition of $A$ defines an equivalence
relation. This fundamental observation will reappear when we study quotient
vector spaces in 
ef{ch:linear-maps}.
\end{remark}

\subsection{Order relations}\label{subsec:order}

\begin{definition}{Partial order}{def:partial-order}
A binary relation $\leq$ on $A$ is a \emph{partial
order}\index{relation!order} if it is:
\begin{enumerate}[label=(\roman*)]
  \item \textbf{Reflexive:} $a\leq a$ for all $a\in A$.
  \item \textbf{Antisymmetric:} $a\leq b$ and $b\leq a$ imply $a = b$.
  \item \textbf{Transitive:} $a\leq b$ and $b\leq c$ imply $a\leq c$.
\end{enumerate}
If in addition every two elements are comparable ($a\leq b$ or $b\leq a$
for all $a,b\in A$), then $\leq$ is a \emph{total order}\index{total order}.
\end{definition}

\begin{example}{Orders on familiar sets}{ex:orders}
\begin{enumerate}[label=(\alph*)]
  \item $(\R,\leq)$ is totally ordered.
  \item $(\mathcal{P}(S),\subseteq)$ is partially ordered for any set $S$,
    but not totally ordered when $\card{S}\geq 2$ (e.g.\ $\set{1}$ and
    $\set{2}$ are incomparable in $\mathcal{P}(\set{1,2})$).
  \item Divisibility on $\N^*$: $a\mid b$ defines a partial order.
\end{enumerate}
\end{example}


% ============================================================
\section{Algebraic structures}\label{sec:algebraic-structures}
% ============================================================

In this section we introduce the hierarchy \emph{group} $\to$
\emph{ring} $\to$ \emph{field}. The central notion for this course is
that of a \emph{field}, since linear algebra is the study of vector
spaces \emph{over a field}. However, groups and rings appear naturally
in the theory (symmetry groups, polynomial rings, matrix rings), so a
brief treatment is warranted.

% ----- Groups -----
\subsection{Groups}\label{subsec:groups}

\begin{definition}{Group}{def:group}
A \emph{group}\index{group} is a pair $(G,\star)$ where $G$ is a nonempty
set and $\star\colon G\times G\to G$ is a binary operation satisfying:
\begin{enumerate}[label=\textbf{G\arabic*.}]
  \item \textbf{Associativity:} $(a\star b)\star c = a\star(b\star c)$
    for all $a,b,c\in G$.
  \item \textbf{Identity:} there exists $e\in G$ such that
    $e\star a = a\star e = a$ for all $a\in G$.
  \item \textbf{Inverses:} for each $a\in G$ there exists $a'\in G$ such
    that $a\star a' = a'\star a = e$.
\end{enumerate}
If in addition $a\star b = b\star a$ for all $a,b\in G$, the group is
called \emph{abelian}\index{group!abelian} (or \emph{commutative}).
\end{definition}

\begin{example}{Basic examples of groups}{ex:groups}
\begin{enumerate}[label=(\alph*)]
  \item $(\Z,+)$, $(\Q,+)$, $(\R,+)$, $(\C,+)$ are abelian groups with
    identity $0$ and inverse $-a$.
  \item $(\Q^*,\times)$, $(\R^*,\times)$, $(\C^*,\times)$ are abelian
    groups with identity $1$ and inverse $1/a$. (Here
    $\Q^*=\Q\setminus\set{0}$, etc.)
  \item $(\GL_n(\R),\times)$, the set of invertible $n\times n$ real
    matrices under matrix multiplication, is a group (non-abelian for
    $n\geq 2$). This group will play a central role starting from
    
ef{ch:determinants}.
  \item $(\Z/n\Z,+)$ is an abelian group of order~$n$.
\end{enumerate}
\end{example}

\begin{proposition}{Uniqueness of identity and inverses}{prop:unique-id}
In any group $(G,\star)$:
\begin{enumerate}[label=(\roman*)]
  \item The identity element is unique.
  \item Each element has a unique inverse.
\end{enumerate}
\end{proposition}

\begin{proof}
(i)~Suppose $e$ and $e'$ are both identities. Then $e = e\star e' = e'$.

(ii)~Suppose $a'$ and $a''$ are both inverses of~$a$. Then
$a' = a'\star e = a'\star(a\star a'') = (a'\star a)\star a'' = e\star a'' = a''$.
\end{proof}

% ----- Rings -----
\subsection{Rings}\label{subsec:rings}

\begin{definition}{Ring}{def:ring}
A \emph{ring}\index{ring} is a triple $(R,+,\cdot\,)$ where:
\begin{enumerate}[label=\textbf{R\arabic*.}]
  \item $(R,+)$ is an abelian group (with identity $0_R$).
  \item Multiplication is associative:
    $(a\cdot b)\cdot c = a\cdot(b\cdot c)$.
  \item There exists a multiplicative identity $1_R$ with
    $1_R\cdot a = a\cdot 1_R = a$ for all~$a$.
  \item Distributivity:
    $a\cdot(b+c) = a\cdot b + a\cdot c$ and
    $(a+b)\cdot c = a\cdot c + b\cdot c$.
\end{enumerate}
If $a\cdot b = b\cdot a$ for all $a,b\in R$, the ring is
\emph{commutative}\index{ring!commutative}.
\end{definition}

\begin{example}{Rings}{ex:rings}
\begin{enumerate}[label=(\alph*)]
  \item $(\Z,+,\times)$ is a commutative ring.
  \item The polynomial ring $\R[X]$ is a commutative ring.
  \item $(\Mat_n(\R),+,\times)$, the set of $n\times n$ real matrices, is
    a ring that is \emph{not} commutative for $n\geq 2$.
  \item $(\Z/n\Z,+,\times)$ is a commutative ring.
\end{enumerate}
\end{example}

% ----- Fields -----
\subsection{Fields}\label{subsec:fields}

The field axioms are the single most important algebraic prerequisite for
this course. A field is, roughly speaking, a set where one can add,
subtract, multiply, and divide (except by zero) and where all the usual
algebraic rules hold.

\begin{definition}{Field}{def:field}
A \emph{field}\index{field} is a triple $(\K,+,\cdot\,)$ satisfying the
following axioms. Let $a,b,c$ denote arbitrary elements of~$\K$.

\medskip
\noindent\textbf{Addition axioms:}
\begin{enumerate}[label=\textbf{A\arabic*.}]
  \item \textbf{Closure:} $a+b\in\K$.
  \item \textbf{Associativity:} $(a+b)+c = a+(b+c)$.
  \item \textbf{Commutativity:} $a+b = b+a$.
  \item \textbf{Identity:} there exists $0\in\K$ with $a+0=a$.
  \item \textbf{Inverses:} for each $a$ there exists $-a\in\K$ with
    $a+(-a)=0$.
\end{enumerate}

\medskip
\noindent\textbf{Multiplication axioms:}
\begin{enumerate}[label=\textbf{M\arabic*.}]
  \item \textbf{Closure:} $a\cdot b\in\K$.
  \item \textbf{Associativity:} $(a\cdot b)\cdot c = a\cdot(b\cdot c)$.
  \item \textbf{Commutativity:} $a\cdot b = b\cdot a$.
  \item \textbf{Identity:} there exists $1\in\K$, $1\neq 0$, with
    $a\cdot 1 = a$.
  \item \textbf{Inverses:} for each $a\neq 0$ there exists $a^{-1}\in\K$
    with $a\cdot a^{-1}=1$.
\end{enumerate}

\medskip
\noindent\textbf{Distributive law:}
\begin{enumerate}[label=\textbf{D\arabic*.}]
  \item $a\cdot(b+c)=a\cdot b + a\cdot c$.
\end{enumerate}
\end{definition}

\begin{remark}{Field = commutative division ring}{rem:field-division-ring}
Equivalently, a field is a commutative ring in which every nonzero element
has a multiplicative inverse. The requirement $1\neq 0$ ensures that the
field has at least two elements.
\end{remark}

\begin{proposition}{Elementary consequences of the field axioms}{prop:field-consequences}
Let $\K$ be a field and $a,b\in\K$. Then:
\begin{enumerate}[label=(\roman*)]
  \item $a\cdot 0 = 0$.
  \item $(-1)\cdot a = -a$.
  \item $a\cdot b = 0$ implies $a=0$ or $b=0$
    (a field has no zero divisors).
\end{enumerate}
\end{proposition}

\begin{proof}
(i)~$a\cdot 0 = a\cdot(0+0)=a\cdot 0 + a\cdot 0$. Adding $-(a\cdot 0)$ to
both sides gives $0 = a\cdot 0$.

(ii)~$a + (-1)\cdot a = 1\cdot a + (-1)\cdot a = (1+(-1))\cdot a = 0\cdot a
= 0$, so $(-1)\cdot a$ is the additive inverse of~$a$.

(iii)~If $a\neq 0$, multiply both sides of $ab=0$ by $a^{-1}$:
$b = 1\cdot b = (a^{-1}a)b = a^{-1}(ab) = a^{-1}\cdot 0 = 0$.
\end{proof}

We now examine the three families of fields that appear most frequently in
linear algebra.

% --- The real numbers ---
\begin{example}{The real numbers $\R$}{ex:reals}
$(\R,+,\times)$ is the prototypical field, and the default setting for
most of undergraduate linear algebra. In addition to the field axioms,
$\R$ carries a total order compatible with the field operations and
satisfies the \emph{completeness axiom} (every nonempty bounded-above
subset has a least upper bound). These extra properties are not required
by the field axioms and are not shared by all fields.
\end{example}

% --- The complex numbers ---
\begin{example}{The complex numbers $\C$}{ex:complex}
Define $\C \deff \setcond{a+bi}{a,b\in\R}$ where $i^2=-1$. Addition and
multiplication are given by:
\begin{align*}
  (a+bi)+(c+di) &= (a+c)+(b+d)i, \\
  (a+bi)(c+di)  &= (ac-bd)+(ad+bc)i.
\end{align*}
The additive identity is $0$, the multiplicative identity is $1$, and the
inverse of $a+bi\neq 0$ is
\[
  (a+bi)^{-1} = \frac{a}{a^2+b^2} - \frac{b}{a^2+b^2}\,i.
\]
A key advantage of $\C$ over $\R$ is that every polynomial of degree
$n\geq 1$ has exactly $n$ roots (counted with multiplicity) --- this is
the \emph{Fundamental Theorem of Algebra}. As a consequence, eigenvalue
theory is cleaner over $\C$, as we shall see in

ef{ch:eigenvalues,ch:spectral}.
\end{example}

% --- Finite fields ---
\begin{example}{Finite fields $\F_p$}{ex:finite-field}
Let $p$ be a prime number. On $\Z/p\Z = \set{[0],[1],\ldots,[p-1]}$,
define addition and multiplication by:
\[
  [a]+[b]\deff[a+b],\qquad [a]\cdot[b]\deff[a\cdot b],
\]
where the operations on the right are performed in $\Z$ and the result is
reduced modulo~$p$. These operations are well-defined (independent of the
choice of representatives) and make $\Z/p\Z$ into a field, denoted
$\F_p$\index{finite field}.

The key point is that $p$ being prime guarantees the existence of
multiplicative inverses: if $[a]\neq[0]$, then $\gcd(a,p)=1$, so by
B\'{e}zout's identity there exist $u,v\in\Z$ with $au+pv=1$, and hence
$[a]\cdot[u]=[1]$.

\medskip\noindent
\textbf{Multiplication table of $\F_5$:}

\smallskip
\begin{center}
\begin{tabular}{c|ccccc}
$\times$ & $[0]$ & $[1]$ & $[2]$ & $[3]$ & $[4]$ \\
\midrule
$[0]$ & $[0]$ & $[0]$ & $[0]$ & $[0]$ & $[0]$ \\
$[1]$ & $[0]$ & $[1]$ & $[2]$ & $[3]$ & $[4]$ \\
$[2]$ & $[0]$ & $[2]$ & $[4]$ & $[1]$ & $[3]$ \\
$[3]$ & $[0]$ & $[3]$ & $[1]$ & $[4]$ & $[2]$ \\
$[4]$ & $[0]$ & $[4]$ & $[3]$ & $[2]$ & $[1]$ \\
\end{tabular}
\end{center}

\smallskip
From the table we read off: $[2]^{-1}=[3]$, $[3]^{-1}=[2]$,
$[4]^{-1}=[4]$.
\end{example}

\begin{remark}{Why primality matters}{rem:why-prime}
If $n$ is not prime, $\Z/n\Z$ is \emph{not} a field. For instance, in
$\Z/6\Z$ we have $[2]\cdot[3]=[0]$ with $[2]\neq[0]$ and $[3]\neq[0]$,
violating the no-zero-divisors property. More generally, finite fields of
order $q$ exist if and only if $q$ is a prime power $p^k$; the
construction for $k\geq 2$ requires quotient rings of polynomial rings and
is beyond our scope.
\end{remark}

\begin{definition}{Characteristic of a field}{def:characteristic}
The \emph{characteristic}\index{characteristic} of a field $\K$, denoted
$\Char(\K)$, is the smallest positive integer $p$ such that
$\underbrace{1+1+\cdots+1}_{p} = 0$ in $\K$. If no such integer exists,
we set $\Char(\K)=0$.
\end{definition}

\begin{example}{Characteristics}{ex:char}
$\Char(\Q) = \Char(\R) = \Char(\C) = 0$, while $\Char(\F_p) = p$.
Throughout this course, \emph{unless otherwise stated, $\K$ denotes a
field of characteristic zero} (typically $\R$ or $\C$).
\end{example}

\begin{remark}{Convention: the symbol $\K$}{rem:convention-K}
In these notes, the letter $\K$ always denotes the base field. Most
results in Chapters~1--5 hold over an arbitrary field; we shall
explicitly note when additional hypotheses on $\K$ (such as
$\Char(\K)=0$, or algebraic closure) are needed.
\end{remark}


% ============================================================
\section{Proof techniques}\label{sec:proofs}
% ============================================================

Proofs are the backbone of mathematics. In this section we review the
principal strategies. The reader should view these not as abstract logical
rules but as practical tools that will be used repeatedly throughout the
course.

\subsection{Direct proof}\label{subsec:direct}

To prove ``if $P$ then $Q$'', assume $P$ is true and deduce $Q$ through a
chain of logical steps.

\begin{example}{A direct proof}{ex:direct-proof}
\emph{Claim:} If $n$ is an even integer, then $n^2$ is even.

\emph{Proof:} Assume $n$ is even. Then $n=2k$ for some $k\in\Z$. Hence
$n^2 = 4k^2 = 2(2k^2)$, which is even. \qed
\end{example}

\subsection{Proof by contrapositive}\label{subsec:contrapositive}

The statement ``if $P$ then $Q$'' is logically equivalent to ``if
$\lnot Q$ then $\lnot P$''. Sometimes the contrapositive is easier to
prove.

\begin{example}{A contrapositive proof}{ex:contrapositive-proof}
\emph{Claim:} If $n^2$ is odd, then $n$ is odd.

\emph{Proof (contrapositive):} We prove: if $n$ is even, then $n^2$ is
even. This is exactly the statement proved in

ef{ex:ex:direct-proof}. \qed
\end{example}

\subsection{Proof by contradiction}\label{subsec:contradiction}

To prove a statement $P$, assume $\lnot P$ and derive a logical
contradiction.

\begin{example}{Irrationality of $\sqrt{2}$}{ex:sqrt2}
\emph{Claim:} $\sqrt{2}\notin\Q$.

\emph{Proof:} Suppose for contradiction that $\sqrt{2}=a/b$ with
$a,b\in\Z$, $b\neq 0$, and $\gcd(a,b)=1$. Then $2b^2=a^2$, so $a^2$ is
even, hence $a$ is even: $a=2k$. Then $2b^2=4k^2$, so $b^2=2k^2$, hence
$b$ is even. But then $\gcd(a,b)\geq 2$, contradicting
$\gcd(a,b)=1$. \qed
\end{example}

\subsection{Proof by induction}\label{subsec:induction}

\begin{definition}{Principle of mathematical induction}{def:induction}
Let $P(n)$ be a statement depending on $n\in\N$. If
\begin{enumerate}[label=(\roman*)]
  \item \textbf{Base case:} $P(n_0)$ is true for some $n_0\in\N$, and
  \item \textbf{Inductive step:} for every $n\geq n_0$, $P(n)$ implies
    $P(n+1)$,
\end{enumerate}
then $P(n)$ is true for all $n\geq n_0$.
\end{definition}

\begin{remark}{Strong induction}{rem:strong-induction}
In the \emph{strong} form of induction, the inductive hypothesis in step
(ii) is strengthened: one assumes $P(k)$ for all $n_0\leq k\leq n$ and
deduces $P(n+1)$. The two forms are logically equivalent.
\end{remark}

\begin{example}{Sum of the first $n$ integers}{ex:sum-integers}
\emph{Claim:} For all $n\geq 1$,
$\displaystyle\sum_{k=1}^{n}k = \frac{n(n+1)}{2}$.

\emph{Proof by induction.}

\textbf{Base case ($n=1$):}
$\displaystyle\sum_{k=1}^{1}k = 1 = \frac{1\cdot 2}{2}$. \checkmark

\textbf{Inductive step:} Assume the formula holds for some $n\geq 1$.
Then
\[
  \sum_{k=1}^{n+1}k = \left(\sum_{k=1}^{n}k\right) + (n+1)
  = \frac{n(n+1)}{2} + (n+1)
  = \frac{n(n+1)+2(n+1)}{2}
  = \frac{(n+1)(n+2)}{2}.
\]
This is the formula with $n$ replaced by $n+1$. \qed
\end{example}

\begin{example}{A preview: dimension of a subspace}{ex:induction-linalg}
Induction is the primary tool for proving results about finite-dimensional
vector spaces. For instance, one proves by induction on $\dim V$ that
every subspace of a finite-dimensional vector space is itself
finite-dimensional (see 
ef{ch:vector-spaces}). The base case is
$\dim V = 0$ (the trivial space), and the inductive step reduces a space
of dimension $n+1$ to one of dimension~$n$.
\end{example}

\subsection{Proof by double inclusion}\label{subsec:double-inclusion}

To show that two sets $A$ and $B$ are equal, one proves $A\subseteq B$
and $B\subseteq A$ separately. This technique is used constantly in
linear algebra (e.g.\ showing that $\Ker f = \set{\vzero}$).

\subsection{Proof by double counting / dimension argument}\label{subsec:dim-argument}

Later in the course, a powerful and characteristically linear-algebraic
proof technique will emerge: the \emph{dimension argument}. To show that
two subspaces $U$ and $W$ are equal, one often shows $U\subseteq W$ (or
$W\subseteq U$) and $\dim U = \dim W$. This will be formalized in

ef{ch:vector-spaces}.


% ============================================================
\section{Exercises}\label{sec:ch0-exercises}
% ============================================================

\begin{exercise}{Set operations \basic}{exo:set-ops}
Let $A=\set{1,2,3,4,5}$ and $B=\set{3,4,5,6,7}$.
Compute $A\cup B$, $A\cap B$, $A\setminus B$, $B\setminus A$, and
$A\times\set{0,1}$.
\end{exercise}

\begin{exercise}{Injectivity and surjectivity \basic}{exo:inj-surj}
For each of the following mappings, determine whether it is injective,
surjective, bijective, or none.
\begin{enumerate}[label=(\alph*)]
  \item $f\colon\Z\to\Z$, $f(n)=2n$.
  \item $g\colon\Z\to\Z$, $g(n)=n+3$.
  \item $h\colon\R\to\R$, $h(x)=x^3-x$.
  \item $\vphi\colon\R\to\R_{>0}$, $\vphi(x)=e^x$.
\end{enumerate}
\end{exercise}

\begin{exercise}{Composition and invertibility \intermediate}{exo:comp-inj-surj}
Prove parts (ii)--(iv) of 
ef{prop:prop:comp-inj-surj}.
\end{exercise}

\begin{exercise}{Equivalence relations \basic}{exo:equiv-rel}
On $\Z\times(\Z\setminus\set{0})$, define $(a,b)\sim(c,d)$ if and only if
$ad=bc$. Verify that $\sim$ is an equivalence relation. What familiar set
does the quotient $(\Z\times(\Z\setminus\set{0}))/{\sim}$ represent?
\end{exercise}

\begin{exercise}{Verifying group axioms \basic}{exo:group-axioms}
Verify that $(\Z/5\Z\setminus\set{[0]},\times)$ is an abelian group. What
is the inverse of $[3]$?
\end{exercise}

\begin{exercise}{Why $\Z$ is not a field \basic}{exo:Z-not-field}
Which field axiom fails for $(\Z,+,\times)$? Give a specific
counterexample.
\end{exercise}

\begin{exercise}{Arithmetic in $\F_7$ \intermediate}{exo:F7}
\begin{enumerate}[label=(\alph*)]
  \item Write out the complete multiplication table of $\F_7$.
  \item Find $[3]^{-1}$, $[5]^{-1}$, and $[6]^{-1}$ in $\F_7$.
  \item Solve the equation $[3]x + [2] = [5]$ in $\F_7$.
\end{enumerate}
\end{exercise}

\begin{exercise}{Complex arithmetic \basic}{exo:complex-arith}
Let $z=3+4i$ and $w=1-2i$. Compute $z+w$, $zw$, $\abs{z}$, $z^{-1}$, and
$z/w$.
\end{exercise}

\begin{exercise}{Field of four elements \challenging}{exo:F4}
We know $\F_p$ exists for every prime $p$. A field of order $4$ also
exists, but it is \emph{not} $\Z/4\Z$ (which has zero divisors). Consider
the set $\F_4=\set{0,1,\alpha,\alpha+1}$ with the rule $\alpha^2+\alpha+1=0$
(i.e.\ $\alpha^2=\alpha+1$) and characteristic~$2$ (so $1+1=0$, and
$x=-x$ for all~$x$).
\begin{enumerate}[label=(\alph*)]
  \item Write out the complete addition and multiplication tables for
    $\F_4$.
  \item Verify that $\F_4$ is indeed a field.
\end{enumerate}
\end{exercise}

\begin{exercise}{Induction: sum of squares \intermediate}{exo:sum-squares}
Prove by induction that for all $n\geq 1$:
\[
  \sum_{k=1}^{n}k^2 = \frac{n(n+1)(2n+1)}{6}.
\]
\end{exercise}

\begin{exercise}{Induction: divisibility \intermediate}{exo:induction-div}
Prove by induction that $3^{2n}-1$ is divisible by $8$ for all $n\geq 1$.
\end{exercise}

\begin{exercise}{Proof strategies \challenging}{exo:proof-strategies}
Prove each of the following statements using the most appropriate proof
technique. State which technique you use.
\begin{enumerate}[label=(\alph*)]
  \item If $a,b\in\R$ and $a+b$ is irrational, then at least one of
    $a,b$ is irrational.
  \item For every $n\in\N$, $n^2+n$ is even.
  \item There is no smallest positive rational number.
  \item For all $n\geq 1$, $\displaystyle\sum_{k=0}^{n}2^k = 2^{n+1}-1$.
\end{enumerate}
\end{exercise}


% ============================================================
\section*{Chapter summary}\label{sec:ch0-summary}
\addcontentsline{toc}{section}{Chapter summary}
% ============================================================

\begin{itemize}
  \item A \textbf{set} is a collection of distinct objects. The standard
    operations are union, intersection, difference, and Cartesian product.
    The number sets $\N\subset\Z\subset\Q\subset\R\subset\C$ form the
    foundation for all that follows.

  \item A \textbf{mapping} $f\colon A\to B$ assigns to each element of
    $A$ exactly one element of $B$. The fundamental classification is into
    injective, surjective, and bijective maps.

  \item An \textbf{equivalence relation} partitions a set into
    equivalence classes. An \textbf{order relation} provides a way to
    compare elements.

  \item A \textbf{group} $(G,\star)$ captures the idea of a reversible
    operation. A \textbf{ring} $(R,+,\cdot\,)$ has two compatible
    operations. A \textbf{field} $(\K,+,\cdot\,)$ is a commutative ring
    in which every nonzero element has a multiplicative inverse.

  \item The principal fields in this course are $\R$ (the reals), $\C$
    (the complex numbers), and $\F_p$ (finite fields of prime order).

  \item The main proof techniques are: direct proof, contrapositive,
    contradiction, mathematical induction, and double inclusion. Mastering
    these is essential for the remainder of the course.
\end{itemize}
